\section{\texttt{env.ref()} and XML IDs}

\begin{frame}[fragile]{What Is \texttt{env.ref()}?}
	\begin{itemize}
		\item \textbf{\texttt{env.ref()}} finds a record by its XML / external ID.
\begin{block}{Method Syntax}
\begin{lstlisting}
record = self.env.ref('module_name.xml_id')
record = self.env.ref('module_name.xml_id', raise_if_not_found=False)
\end{lstlisting}
\end{block}
		\item Extremely common for groups, stages, sequences, default records, and actions.
	\end{itemize}
\end{frame}

\begin{frame}[fragile]{Practical Examples}
\begin{lstlisting}
admin_group = self.env.ref('base.group_system')
student_action = self.env.ref('school_manager.action_student')
draft_stage = self.env.ref('project.project_stage_0', raise_if_not_found=False)

if self.env.user.has_group('base.group_system'):
	# admin-only logic
	pass
\end{lstlisting}
\end{frame}

\begin{frame}[fragile]{Why Prefer XML IDs over Hard-Coded IDs?}
	\begin{itemize}
		\item Database IDs change between databases.
		\item XML IDs are stable across installs when defined in data files.
\begin{lstlisting}
# Bad
stage = self.env['project.task.type'].browse(3)

# Good
stage = self.env.ref('project.project_stage_0')
\end{lstlisting}
		\item Use \texttt{raise\_if\_not\_found=False} when the record may be optional.
	\end{itemize}
\end{frame}

\begin{frame}[fragile]{Other Useful Environment Helpers}
\begin{lstlisting}
self.env.lang          # current language, e.g. 'en_US'
self.env.is_superuser()
self.env.is_admin()
self.env.is_system()
\end{lstlisting}
	\begin{itemize}
		\item These helpers make environment checks clearer and safer.
	\end{itemize}
\end{frame}
